\scanpage{189}
$X/S$ \quad schéma sur $S$.

1) Un module \inserted{qu. coh.} $E$ sur $X$, avec connexion rel/$S$, est dit \deleted{\unsure{le|de}} \uline{trivial} si $E\simeq f^\ast(F)$, $F$ un module \inserted{qu. coh.} sur $S$ [isom respectant connexion]. \margin{Il serait plus naturel à certains égards de travailler avec les PD-connexions [c'est pareil dans cas lisse]} \uline{NB} Si $X\longrightarrow S$ est \deleted{\unsure{couvrant fppf|connexe fpqc}} \inserted{lisse surjectif} et à fibres géométriquement connexes, \inserted{$S$ de carac 0, alors,} le foncteur $F\longrightarrow f^\ast(F)$, des modules qu. coh. sur $S$ dans les modules qu. coh. sur $X$ à connexion rel/$S$, est pleinement fidèle [du moins si $F$ de présentation finie].

On dit que $E$ est \uline{loc. trivial} ou \uline{Zar-intégrable} si $E$ est trivial localm\textsuperscript{t} sur $(X\longrightarrow S)$ pour la topologie de Zariski. On dit que $E$ est \uline{intégrable}, si la m\textsuperscript{e} condition est satisfaite, loc. pour top. étale.

2) Supposons $S$ de car $p>0$, et que $X$ sur $S$ admette loc. une $p$-base. Alors \deleted{$E$ est loc. trivial} conditions suivantes sont équivalentes sur $E$ qu coh.

a) $E$ est \uline{Zar. intégrable}

b) $E$ est intégrable

c) $E$ est à courbure et $p$-courbure nulle.

d) $E$ est \uline{trivial} relativement à $X\longrightarrow(X^{p/S},\underline O_{X^{p/S}}/J\simeq\Lambda\deleted{\unsure{\underline O_S|{}_{\!S}}})$

De plus, le foncteur $F\mathrel{\text{---}}{f'}_{X/S}^\ast(F)$ des Modules cohérents sur $X^{p/S\,\prime}=(X^{p/S},\underline O_{X^{p/S}}/J=\Lambda)$ aux modules \inserted{qu. coh. sur $X$} à connexion intégrable rel/$S$, est une équivalence de catégories.
% archive sheet number: 188
